\usepackage[utf8]{inputenc}
\usepackage[T1]{fontenc}
\usepackage[czech]{babel}
\usepackage{lmodern}
\usepackage[margin=2.2cm]{geometry}
\usepackage{amsmath,amssymb,amsthm,mathtools}
\usepackage{physics}
\usepackage[dvipsnames]{xcolor}
\usepackage{tikz}
\usetikzlibrary{arrows.meta,calc,decorations.pathmorphing,decorations.pathreplacing,patterns,positioning,babel}
\usepackage{tikz-cd}
\usepackage{enumitem}
\usepackage[normalem]{ulem}
\usepackage[skins,breakable,theorems]{tcolorbox}
\usepackage{titlesec}
\usepackage[colorlinks=true,linkcolor=NavyBlue,urlcolor=NavyBlue]{hyperref}

%======================================================================
% Barvy
%======================================================================
\definecolor{hlavni}{RGB}{0,70,140}      % nadpisy
\definecolor{defbarva}{RGB}{30,100,200}  % definice
\definecolor{vetabarva}{RGB}{190,40,40}  % věty, postuláty
\definecolor{prbarva}{RGB}{20,130,70}    % příklady
\definecolor{vzorbarva}{RGB}{220,120,0}  % důležité vzorce

\titleformat{\section}{\color{hlavni}\Large\bfseries}{\thesection}{1em}{}[{\color{hlavni}\titlerule[1pt]}]
\titleformat{\subsection}{\color{hlavni}\large\bfseries}{\thesubsection}{1em}{}
\titleformat{\subsubsection}{\color{hlavni!80!black}\normalsize\bfseries}{\thesubsubsection}{1em}{}

%======================================================================
% Zvýraznění
%======================================================================
% Definice (modrý box)
\newtcolorbox{definice}[1][]{enhanced,breakable,colback=defbarva!6,colframe=defbarva,
  boxrule=0.6pt,left=4pt,right=4pt,top=2pt,bottom=2pt,arc=2pt,
  fonttitle=\bfseries\small,coltitle=white,title={Definice\ifx&#1&\else: #1\fi}}
% Věta / postulát / lemma / klíčové tvrzení (červený box)
\newtcolorbox{veta}[1][Věta]{enhanced,breakable,colback=vetabarva!5,colframe=vetabarva,
  boxrule=0.6pt,left=4pt,right=4pt,top=2pt,bottom=2pt,arc=2pt,
  fonttitle=\bfseries\small,coltitle=white,title={#1}}
% Příklad (zelený proužek vlevo)
\newtcolorbox{priklad}[1][Příklad]{enhanced,breakable,colback=prbarva!4,colframe=prbarva,
  boxrule=0pt,leftrule=2.5pt,left=5pt,right=3pt,top=2pt,bottom=2pt,sharp corners,
  fonttitle=\bfseries\small,coltitle=prbarva!80!black,colbacktitle=prbarva!12,title={#1}}
% Vysvětlení doplněné při přepisu (není v původních poznámkách) – fialový box
\definecolor{vysvbarva}{RGB}{110,60,160}
\newtcolorbox{vysvetleni}[1][]{enhanced,breakable,colback=vysvbarva!5,colframe=vysvbarva,
  boxrule=0.5pt,leftrule=2.5pt,left=5pt,right=4pt,top=2pt,bottom=2pt,arc=1pt,
  fontupper=\small,fonttitle=\bfseries\small,coltitle=white,colbacktitle=vysvbarva,
  title={Vysvětlení\ifx&#1&\else: #1\fi}}
% Oprava zjevné chyby v poznámkách: \oprava{správně}{v originále}
\definecolor{opravbarva}{RGB}{0,128,128}
\newcommand{\oprava}[2]{{\color{opravbarva}#1}\relax\ifmmode\text{\color{opravbarva}\scriptsize\itshape\,[oprava; orig.: $#2$]}\else{\color{opravbarva}\scriptsize\itshape\,[oprava; orig.: #2]}\fi}
% Důležitý vzorec v rámečku (v matematice):  \vzorec{E = mc^2}
\newcommand{\vzorec}[1]{\tcbhighmath[colback=vzorbarva!8,colframe=vzorbarva,boxrule=0.8pt,arc=2pt,left=3pt,right=3pt,top=2pt,bottom=2pt]{#1}}
% Klíčový pojem v textu
\newcommand{\kw}[1]{\textbf{\color{hlavni}#1}}
% Upozornění / pozor
\newcommand{\pozor}[1]{{\color{vetabarva}\bfseries #1}}

% Poznámky psané tužkou (v originále slabě, šedě)
\newcommand{\pozn}[1]{\relax\ifmmode\text{\color{gray}\small\itshape [#1]}\else{\color{gray}\small\itshape [#1]}\fi}
% Číslo strany originálních poznámek
\newcommand{\strana}[1]{\marginpar{\colorbox{hlavni!10}{\small\texttt{\color{hlavni}#1}}}}

\setlength{\parindent}{0pt}
\setlength{\parskip}{4pt}
